\chapter{APÊNDICE}
\label{ap:codigo}

O equacionamento do robô foi implementado em Python, no ambiente Jupyter, no notebook \texttt{blocos\_inercia\_robo.ipynb}, organizado em onze blocos cumulativos. A seguir, a estrutura (\ref{tab:blocos}):

\begin{table}[H]
\centering
\caption{\label{tab:blocos}- Estrutura do notebook \texttt{blocos\_inercia\_robo.ipynb}.}
\vspace*{.1cm}
\small
\begin{tabular}{|c|p{11.5cm}|}
\hline
\textbf{Bloco} & \textbf{Conteúdo} \\
\hline
1 & Importações (NumPy, SymPy, Matplotlib) e configuração de exibição. \\
\hline
2 & Funções gerais: exibição de equações em LaTeX, conversão de unidades e funções utilitárias. \\
\hline
3 & Braço: massas, centro de massa e tensor de inércia das três barras, com cálculo simbólico em SymPy. \\
\hline
4 & Perna (contrapeso): mesma rotina do Braço, com desenho cotado da peça. \\
\hline
5 & Ombro: manivela em pilha de camadas, furo passante (massa negativa) e pino cilíndrico. \\
\hline
6 & Motor com caixa de redução: composição do corpo e da caixa a partir do desenho técnico. \\
\hline
7 & Bateria: paralelepípedo homogêneo, verificação de densidade e conversão de unidades. \\
\hline
8 & Quadril: haste fina, comparação entre as hipóteses de seção circular e quadrada. \\
\hline
9 & Montagem do robô completo no sistema do corpo, com a posição ($\psi=0^\circ$) definida na Tabela de dimensões e o desenho das vistas. \\
\hline
10 & Corpos móveis e manivela a 30 rpm: cinemática das alavancas (Equações \ref{eq:phi} e \ref{eq:dphi}) e inércia em torno dos pivôs. \\
\hline
11 & Modelo dinâmico plano: montagem simbólica das matrizes de massa e dos vetores de contato, verificação contra Lagrange, detecção de eventos de engate e integração das equações da marcha na corda. \\
\hline
\end{tabular}
\end{table}
